% No preamble. Prefix cgd-. Depends on the checked covariant tangent bridge.
\section{Actual local densities in the covariant mean-zero gauge}
\label{cgd-section}

Use the covariant tangent construction above, with all-site
fine gauge generator \(D\), covariant scalar block average
\(B\), corrected bond differential \(\bar C\), and
\[
 BB^T=K^{-4}I,\qquad
 P=I-K^4B^TB,\qquad
 \widehat G=P D^TD P+K^{-2}(I-P),\qquad
 Q_B=\mathcal H+D P D^T+a\bar C^T\bar C.
\]
The actual finite determinant identity is
\begin{equation}
 D_{M,K}=\mathcal N_0
       \left(\log\det\widehat G-\tfrac12\log\det Q_B\right).
 \label{cgd-identity}
\end{equation}
Here \(\mathcal H\) is the full residual Hessian, not a
substitution of the bare Wilson Hessian for \(S_M''\).
The matrix \(P\) projects onto scalar fields with zero
covariant block mean. It is not the projection onto
fields zero at the block centers.

\subsection{Gauge covariance and localized logarithms}
\label{cgd-covariance-section}

Let a fine site gauge induce orthogonal adjoint matrices
\(O_S,O_E,O_c\) on fine scalars, fine bonds, and coarse
scalars. The covariant tangent construction gives
\[
 B(U^g)=O_c B(U)O_S^T,\qquad
 P(U^g)=O_S P(U)O_S^T,\qquad
 \bar C(U^g)=O_{E_c}\bar C(U)O_E^T.
\]
Together with the exact covariance of \(D,\mathcal H\)
this proves
\begin{equation}
 \widehat G(U^g)=O_S\widehat G(U)O_S^T,\qquad
 Q_B(U^g)=O_EQ_B(U)O_E^T.
 \label{cgd-covariance}
\end{equation}
These are also complex-orthogonal identities on the
holomorphic germs. No choice of a frame for \(\ker B\)
appears in them.

Fix \(K\). The flat lower bounds for both matrices are
\(cK^{-2}\), uniformly in old depth and volume. Their
flat upper bounds and exponentially localized kernels,
and their first-two source kernels after actual
stationary substitution, are uniform in those same
parameters at this fixed \(K\).
The block-local scalar projection does not enlarge
support beyond the corresponding final block; the
bond construction has bounded range in final units.
The old action and stationary response retain their
established exponential source bounds.

Consequently there is a complex output ball, allowed
to depend on \(K\), on which both logarithms have
source-rooted exponential kernels uniformly in volume
and old depth. One direct proof is the logarithm series.
Choose \(s_0\) larger than the flat upper spectral bound.
At flatness \(\|I-T/s_0\|<1\), for either matrix \(T\).
A sufficiently small exponential conjugation and
background perturbation preserve that strict bound.
Then
\[
 \log T=(\log s_0)I-
       \sum_{n\ge1}\frac{(I-T/s_0)^n}{n}
\]
converges in the resulting weighted matrix algebra.
First and second source derivatives cost at most
polynomial factors in \(n\), absorbed by the geometric
decay. Every source remains attached to the matrix
product chain. Spending a smaller exponential converts
the weighted bounds into absolute rooted sums.
The fine fiber dimension is finite at this fixed \(K\).
No endpoint-independent source constant is inferred
from this argument.

Group scalar sites and positive fine bonds by their
initial final block \(x\), and define
\begin{equation}
 \rho_x(W)=\mathcal N_0\left\{
  \operatorname{Tr}_{x,S}\log\widehat G(W)
       -\tfrac12\operatorname{Tr}_{x,E}\log Q_B(W)\right\}.
 \label{cgd-density}
\end{equation}
The finite sum is \(D_{M,K}\). Each diagonal color
trace is individually gauge invariant by
\eqref{cgd-covariance}. The actual stationary branch
under an output gauge changes by a fine gauge; its
root-gauge convention therefore does not change these
traces. Along the chosen jointly locally convergent structural
system, these densities have local infinite limits
and absolutely summable polynomial source moments,
without introducing a globally based scalar inverse.

In particular, each of the two new factors separately
has zero thermodynamic quadratic response to a
constant commuting connection, at fixed \(K\).
Apply its local Ward identity to a linear scalar
parameter and pass the absolutely summable Hessian
row sum. Thus the growing constant-connection terms
of the globally based physical quotient have
already canceled in the exact change of variables;
they are not discarded by a subsequent limit.

\subsection{Matching the actual Maxwell coefficient}
\label{cgd-Maxwell-section}

\begin{proposition}[Local density and actual affine coefficient]
\label{cgd-density-theorem}
In every joint infinite structural system at fixed \(K\),
the density in \eqref{cgd-density} satisfies
\begin{equation}
 \beta(D_{M,K})=
    2[t^2]\rho_0\bigl(w_\mu(n)=t n_\nu H\bigr),
 \qquad \mu\ne\nu,\quad\tau(H^2)=1.
 \label{cgd-Maxwell}
\end{equation}
The polynomial substitution uses the first two localized
Taylor kernels. It is not an evaluation of a bounded-field
analytic map at a nonzero unbounded connection.
\end{proposition}

\begin{proof}
The linear color jet of each invariant density is zero
for the semisimple color algebra. Its quadratic Ward
identity kills gradients in either slot. Exponential
source moments extend the identity to linear scalar
parameters, so its row and column sums vanish for
each fixed bond component. Write \(k_0(b,c)\) for that
density's Hessian on the chosen component and color.
Then
\[
 [t^2]\rho_0(t n_\nu H)
   =\frac12\sum_{b,c}k_0(b,c)b_\nu c_\nu
   =-\frac14\sum_{b,c}k_0(b,c)(b_\nu-c_\nu)^2.
\]
All sums are absolute. Translation covariance and
reanchoring identify the last expression with
\(-\frac14\sum_r r_\nu^2\mathcal K_{\mu\mu}(0,r)\),
where \(\mathcal K\) is the actual total Hessian
kernel. Its Maxwell moment is
\(-\frac12\sum_r r_\nu^2\mathcal K_{\mu\mu}(0,r)\).
This proves \eqref{cgd-Maxwell}.
\end{proof}

The actual stationary polynomial jets on this probe
are the commuting affine fine connection established
in Notebook 33. All gauge-covariant matrices may
therefore be computed in that representative.
The full effective classical Hessian and the
covariant averaging operators remain in \(Q_B\);
the formula does not replace their curvature
derivatives by a continuum operator.

\subsection{A coupled heat representation with the physical flat scale}
\label{cgd-heat-section}

Let \(\widehat G(t),Q_B(t)\) denote the actual affine
jets just specified, and set
\begin{equation}
 h^{\rm mean}_{M,K}(s)
  =[t^2]\left\{
    \operatorname{Tr}_{0,E}e^{-sQ_B(t)}
             -2\operatorname{Tr}_{0,S}e^{-s\widehat G(t)}
         \right\}.
 \label{cgd-heat-coefficient}
\end{equation}
The fixed-\(K\) logarithm/resolvent identities, or
ordered Duhamel with the localized first-two vertices,
give the exact representation
\begin{equation}
 \boxed{\displaystyle
     \beta(D_{M,K})=
        \int_0^\infty h^{\rm mean}_{M,K}(s)\,\frac{ds}{s}.}
 \label{cgd-heat-identity}
\end{equation}
There is no additional coarse determinant in this
covariant choice: \(BB^T=K^{-4}I\), and the added
scalar complement in \(\widehat G\) contributes only
a background-independent determinant.

For precision, this is a fixed-\(K\) identity on
localized Taylor kernels. The spectral gaps,
exponential weighted kernels and finite polynomial
vertices control the first-two Duhamel coefficients
by an integrable small-time and exponentially
decaying large-time envelope. Its constants may
depend on \(K\). This justifies the logarithm
integral, polynomial source substitution and local
limit; it does not give a uniform curvature
estimate as \(K\) grows.

The new scalar and vector factors both have flat
spectral scale \(K^{-2}\). The point-pinned ghost
coefficient \(4\alpha_g\log K+O(1)\) of Notebook 34
belongs to a different exact determinant splitting
and must not be assigned to \(\widehat G\).
The sum in \eqref{cgd-heat-identity} is the same
actual coefficient. Its endpoint-uniform magnetic
comparison remains to be proved in these new,
local, gauge-covariant variables.
